\section{The system under study}
\label{sec:system}

\subsection{The relation being closed}

The solubility of a crystalline solid in a liquid follows from a balance between two quantities.
Breaking the crystal costs an amount set by the melting temperature $T_m$ and the enthalpy of
fusion $\Delta H_{\mathrm{fus}}$; mixing the freed molecules into the solvent is governed by the
activity coefficient $\gamma_2$ of the solute in that solvent. For an ideal solution the two
combine into the mole fraction at saturation,
%
\begin{equation}
\ln x_2 \;=\; -\,\Phi(T) \;-\; \ln \gamma_2 ,
\qquad
\Phi(T) \;=\; \frac{\Delta H_{\mathrm{fus}}}{R}\!\left(\frac{1}{T} - \frac{1}{T_m}\right),
\label{eq:sle}
\end{equation}
%
with $R$ the gas constant. The first term is the crystal contribution and carries no information
about the solvent; the second is where the chemistry of the mixture enters. Because $\gamma_2$
itself depends on $x_2$, equation~\eqref{eq:sle} is solved rather than evaluated.

COSMO-SAC \cite{lin2002} supplies $\gamma_2$. It represents each molecule by its $\sigma$-profile:
the area of the molecular cavity surface carrying screening-charge density $\sigma$, discretised
onto a fixed grid of $51$ bins spanning roughly $\pm 0.025\,e\,\mathrm{\AA}^{-2}$. A molecule's
identity enters the model through this vector and through the total cavity area; nothing else about
the compound is used. Segments exchange with a common pool, and a self-consistent calculation
returns a segment activity coefficient for each bin, from which the molecular activity coefficient
follows. The exchange energy contains a term that switches on only for segments beyond a threshold
$\sigma_{\mathrm{hb}}$ in magnitude, which is how the model represents hydrogen bonding; a segment
qualifies by its charge density alone, with no reference to which atom carries it. That threshold
divides the charge-density axis into a core and two outer windows, and the division recurs
throughout what follows.

\subsection{The network that supplies the profile}

A graph neural network reads the solute and the solvent and produces, for each, a $\sigma$-profile
and a cavity area; these are passed to the COSMO-SAC calculation and to the solver for
equation~\eqref{eq:sle}. A second head predicts the crystal properties $T_m$ and
$\Delta H_{\mathrm{fus}}$ that enter $\Phi(T)$. The profile head is a feed-forward network on the per-molecule representation, ending in a normalised, non-negative vector of $51$ bins. The same head serves both roles, but reads a role-dependent representation, so a compound receives one profile as solute and another as solvent.

The head is trained in two ways at once, and this is the property that makes the pipeline suitable
for the present question. It is supervised directly against tabulated reference profiles, as a
predictor of those profiles would be. It is also trained through the thermodynamic relation, so
the gradient of the solubility loss reaches it. The nearest prior model instead supervises its
profile predictor externally and then freezes it, so that no gradient from the phase-equilibrium
data reaches the intermediate \cite{yang2025tennet}. The pipeline studied here therefore occupies
the position where both influences act, which is where their relative effect can be read.

\subsection{Reference profiles, data and arms}

Reference $\sigma$-profiles are taken from the VT-2005 tabulation \cite{mullins2006}, computed by
quantum chemistry with a continuum solvation treatment. Throughout, ``reference'' means a value
from that table, and departures are measured against it. How far the tabulation itself depends on
the underlying electronic structure method is a separate question \cite{salih2025sigma} and is not
at issue here, since every comparison in this work uses one table.

Solubility data are drawn from a curated compilation of measured solid solubilities
\cite{krasnov2025bigsoldb}, split so that no scaffold \cite{bemis1996} appears in both training and
test sets, which is the harder of the two standard settings and the one relevant to new compounds.
Infinite-dilution activity coefficients used to score the closure directly come from an independent
compilation. All results below are reported over five training seeds; the arms differ only in the
seed, so the spread across them is the natural scale against which any effect must be read.
Architectural details, hyperparameters, the construction of the splits and the auxiliary
supervision streams are given in the Supporting Information.

\subsection{How the measurements are read}

Each measurement in Section~\ref{sec:results} is a contrast, and each is reported against a control
chosen so that it would fire if the instrument responded to something other than the intended
effect. Where a quantity is compared with a reference table, the control is the same quantity
computed on profiles that carry the same total area but no chemistry. Where part of a profile is
replaced, the control is a part of the same size chosen at random. Where an effect is attributed to
training, the control is an untrained network of the same architecture. The reading rule for each
measurement, together with the result that would refute it, was fixed before the measurement was
made; the rules and their outcomes are listed in the Supporting Information. Table~\ref{tab:battery}
collects every test with the control beside it, because a result of this kind is not interpretable
without one.
